\emph{(i).} The adjoint extends antilinearly to $\Cl(p,q)_{\C}$, and a state satisfies $\omega(y^{*}) = \overline{\omega(y)}$. For orthogonal generators $g, h$ with $g^{2} = h^{2} = s$, $g^{*} = sg$ and $(gh)^{*} = h^{*}g^{*} = hg = -gh$, so $\omega(gh)$ is imaginary, and zero for a real state. With $c = \tfrac12(g + ih)$: $c^{2} = \tfrac14(g^{2} - h^{2}) = 0$, $c^{*} = \tfrac{s}{2}(g - ih)$, $\{c, c^{*}\} = 1$ and $c^{*}c = \tfrac12(1 + is\,gh)$, so $\omega(c^{*}c) = \tfrac12$. For the null space: in the GNS representation with cyclic vector $\xi$, $\omega(y^{*}y) = \|y\xi\|^{2}$, so $N_{\omega} = \{y \in V_{\C} : y\xi = 0\}$ is a subspace, and for $y, z \in N_{\omega}$, $2B(y,z)\xi = (yz + zy)\xi = 0$ (\claimref{std-algebra.vector-squares}), so $N_{\omega}$ is totally isotropic. Complex conjugation $y \mapsto \bar y$ of $V_{\C}$ satisfies $\overline{y}^{*}\,\overline{y} = \overline{y^{*}y}$, and for real $\omega$, $\omega(\bar w) = \overline{\omega(w)}$; so $N_{\omega}$ is conjugation invariant, hence $N_{\omega} = U \otimes \C$ with $U = N_{\omega} \cap V$ real and totally isotropic, of dimension at most the Witt index $\min(p,q)$ (Fact \claimref{tower.witt-index-quantum-capacity}). A Fock vacuum of a maximal isotropic $W$ has $W \subseteq N_{\omega}$ with $\dim W = \lfloor (p+q)/2 \rfloor$, the complex capacity (Proposition \claimref{one-relation.two-capacities-three-consequences}), which exceeds $\min(p,q)$ unless $|p - q| \le 1$. For $(2,4)$ the incidence vacua are vacua of three-pair systems (Fact \claimref{one-relation.imported-pure-spinors-incidence-chart}) and $\min(2,4) = 2 < 3$.

\emph{(ii).} From the definitions, $a^{*} = \nO^{*} = \tfrac12(e_{+} - e_{-})$, $b^{*} = \tfrac12(e_{1} - e_{2})$, and $a, b$ satisfy the CAR (distinct orthogonal null pairs, normalised by $\ip{\nO}{\ninf} = 1$). Since $t e_{1} + r e_{2} = v\,b + u\,b^{*}$ and $\ninf = 2a^{*}$,
\[
X(x) = a + v\,b + u\,b^{*} - uv\,a^{*} = (a + u\,b^{*}) + v\,(b - u\,a^{*}) = (a + v\,b) + u\,(b^{*} - v\,a^{*}) .
\]
Using $a\ket{0} = b\ket{0} = 0$ and the CAR: $(a + u b^{*})(1 - u a^{*}b^{*})\ket{0} = -u\,b^{*}\ket{0} + u\,b^{*}\ket{0} = 0$ and $(b - u a^{*})(1 - u a^{*}b^{*})\ket{0} = u\,a^{*}\ket{0} - u\,a^{*}\ket{0} = 0$, so $X\Omega_{x} = 0$; $(b^{*} - v a^{*})^{2} = 0$ and $(a + v b)(b^{*} - v a^{*})\ket{0} = (-v + v)\ket{0} = 0$, so $X\Psi_{x} = 0$. The element $a^{*}b^{*}$ is even and $b^{*} - v a^{*}$ is odd, so both vectors lie in the half of $\ket{0}_{+}$; they are nonzero and independent (they differ in the parity of the $(a,b)$ occupation), and the kernel of $X$ on that half, the incidence line of Proposition \claimref{one-relation.spacetime-manifold-vacuum-lines}(2), is two-dimensional (Fact \claimref{one-relation.imported-pure-spinors-incidence-chart}), so they span $L_{x}$. They have definite spectator occupation, since $\ket{0}_{+}$ and $\ket{0}_{-}$ differ in the occupation of $c$ and $a, b$ do not act on it. Directly, $\|\Omega_{x}\|^{2} = 1 + u^{2}$ with $\langle N \rangle = u^{2}/(1 + u^{2})$, and $\|\Psi_{x}\|^{2} = 1 + v^{2}$ with $\langle N \rangle = v^{2}/(1 + v^{2})$. The vacua $\ket{0}_{\pm}$ differ in the occupation of the third mode, on which $N$ does not act, so $\Omega_{x} \perp \Psi_{x}$ and $\langle \Omega_{x}, N\Psi_{x}\rangle = 0$. In the $(a, b)$ occupation basis, $\Omega_{x}/\|\Omega_{x}\| = (\ket{00} - u\ket{11})/\sqrt{1 + u^{2}}$ and $\Psi_{x}/\|\Psi_{x}\| = (\ket{01} - v\ket{10})/\sqrt{1 + v^{2}}$, each tensored with a definite spectator state; so the $(a,b)$ state is pure. For $u \ne 0$ the first is in Schmidt form with both coefficients nonzero and $\lambda_{1}^{2}/\lambda_{0}^{2} = u^{2}$, and the ascent marginal $\operatorname{diag}(1, u^{2})/(1 + u^{2})$ is faithful, so by Theorem \claimref{state-selection.partner-purification-are-canonical} it is the thermofield double of that Gibbs state with partner $b$; $e^{-\beta} = u^{2}$ gives $\beta < 0$ exactly when $|u| > 1$. The same holds for $\Psi_{x}$ with $v$, after the unitary $\ket{0} \leftrightarrow \ket{1}$ on the partner. At $u = 0$ the marginal is the pure vacuum, not faithful, and the theorem does not apply. Reality on the subalgebra of $e_{1}, e_{2}, e_{\pm}$ and the value $\pm i$ on $e_{3}e_{4}$: $e_{+} = a + a^{*}$, $e_{-} = a - a^{*}$, $e_{1} = b + b^{*}$, $e_{2} = b - b^{*}$, so every element of that subalgebra is a real combination of words in $a, a^{*}, b, b^{*}$, whose matrix elements between the real vectors $\ket{0}, a^{*}b^{*}\ket{0}, a^{*}\ket{0}, b^{*}\ket{0}$ are real, and $u, v$ are real; $c^{*}c = \tfrac12(1 - i\,e_{3}e_{4})$ with $c$ in a definite occupation gives $\langle e_{3}e_{4} \rangle = \pm i$.

\emph{(iii).} In an orthonormal basis $\hat\Omega_{x}, \hat\Psi_{x}$ of $L_{x}$, $N$ compresses to $\operatorname{diag}(u^{2}/(1+u^{2}), v^{2}/(1+v^{2}))$ (the off-diagonal term vanishes because the spectator occupations differ), so $\langle N \rangle$ on unit vectors of $L_{x}$ takes every value of the segment. Each vacuum on $L_{x}$ has definite chirality and $a$ is odd, so $\langle a \rangle = 0$ and the marginal is diagonal, hence of the descended Gibbs form for nonzero weight (Theorem \claimref{singularity.descended-form-state-independent}, whose zero-weight scope requires occupation $\tfrac12$ there). For $\phi = \lambda\hat\Omega_{x} + \mu\hat\Psi_{x}$ and any word $w$ in $a, a^{*}, b, b^{*}$, the cross terms $\langle \hat\Omega_{x}, w\hat\Psi_{x}\rangle$ vanish because $w$ does not change the spectator occupation; so the $(a,b)$ state of $\phi$ is $|\lambda|^{2}$ times that of $\hat\Omega_{x}$ plus $|\mu|^{2}$ times that of $\hat\Psi_{x}$, a mixture of two orthogonal pure states (of opposite $(a,b)$ parity), mixed when $\lambda\mu \ne 0$. Since $\phi$ is pure on $(a, b, c)$, the purification of the ascent mode is by $b$ and $c$ jointly.

\emph{(iv)} collects (ii) and (iii); the condition for $\beta = 2\pi$ is $e^{-2\pi}$ between $u^{2}$ and $v^{2}$, as $s \mapsto s/(1+s)$ is increasing. Under the dilation $x \mapsto \lambda x$ the null coordinates scale by $\lambda$ while $N$ is fixed, and the implementing rotor is self-adjoint, not unitary, so the segment is not dilation invariant.

\emph{(v).} The plane of $\varepsilon, \nu$ has Gram diagonal $(-1, \nu^{2})$, indefinite exactly when $\nu^{2} = +1$. If $\nu^{2} = -1$, the argument of (i) applied to the two generators $\varepsilon, \nu$ only needs reality of the state on their subalgebra: $\varepsilon\nu$ is skew-adjoint, so its expectation vanishes and the mode $\tfrac12(\varepsilon + i\nu)$ has occupation $\tfrac12$. For $\nu^{2} = +1$, with $n = \tfrac12(\nu + \varepsilon)$, $n^{*}n = \tfrac12(1 + \nu\varepsilon)$ and $\nu\varepsilon$ is self-adjoint, so reality leaves the occupation free. The reality premise itself is not derived.

The check builds the Fock module in exact Gaussian-rational arithmetic and verifies the unitary structure, the CAR, the kernel dimension, $X\Omega_{x} = X\Psi_{x} = 0$, the occupations, the block-diagonal compression, the reality on the four hyperbolic generators and the value $\pm i$ on $e_{3}e_{4}$, at eight rational points; it computes the purity of the $(a,b)$ state from the sixteen Majorana monomials, finding $1$ on $\Omega_{x}$ and $\Psi_{x}$ and a value below $1$ on equal superpositions; and it verifies the identities of (i).
